% from ega2-1-fr.tex
\oldpage[II]{196}
que si $\mathcal{N}$ est déjà un sous-$\mathcal{S}$-Module \emph{gradué} de
$\mathcal{M}$, alors $\overline{\mathcal{N}}$ s'identifie à la somme directe
des composantes $\widehat{\mathcal{N}}_n$ de degré $n\geqslant0$ dans
$\widehat{\mathcal{N}}=\mathcal{N}[z]$.
\end{env}

\begin{proposition}[8.13.3]
\label{II.8.13.3-fr}
Le $\mathcal{O}_{\widehat{C}}$-Module
$\operatorname{Proj}_0(\overline{\mathcal{N}})$ est le prolongement canonique
$(\widetilde{\mathcal{N}})^{-}$ de $\widetilde{\mathcal{N}}$ à
$\widehat{C}$.
\end{proposition}

\begin{proof}
La question est locale sur $Y$ et $\widehat{C}$ en vertu de la définition du
prolongement canonique (\hyperref[I.9.4.1-fr]{I, 9.4.1}). On peut donc déjà
supposer que $Y=\operatorname{Spec}(A)$ est affine, avec
$\mathcal{S}=\widetilde{S}$, $\mathcal{M}=\widetilde{M}$ et
$\mathcal{N}=\widetilde{N}$, où $N$ est un sous-$S$-module non
nécessairement gradué de $M$. En outre
(\hyperref[II.8.3.2.6-fr]{8.3.2.6}) $\widehat{C}$ est réunion des ouverts
affines $\widehat{C}_z=C$ et
$\widehat{C}_f=\operatorname{Spec}(S_f^{\leq})$ ($f$ homogène dans $S_+$).
Il suffira donc de montrer que~: $1^\circ$ la restriction de
$\operatorname{Proj}_0(\overline{\mathcal{N}})$ à $C$ est
$\widetilde{\mathcal{N}}$~; $2^\circ$ la restriction de
$\operatorname{Proj}_0(\overline{\mathcal{N}})$ à chaque $\widehat{C}_f$ est
le prolongement canonique de la restriction de $\mathcal{N}$ à
$C\cap\widehat{C}_f=\operatorname{Spec}(S_f)$
(\hyperref[II.8.3.2.6-fr]{8.3.2.6}). Pour le premier point, notons que
$\operatorname{Proj}_0(\overline{\mathcal{N}})|C$ s'identifie à
$(\overline{N}_{(z)})^{\sim}$
(\hyperref[II.8.3.2.4-fr]{8.3.2.4})~; or $\overline{N}_{(z)}$ s'identifie
canoniquement (\hyperref[II.2.2.5-fr]{2.2.5}) à l'image de
$\overline{N}$ dans $\widehat{M}/(z-1)\widehat{M}$, et par l'isomorphisme
canonique de ce dernier sur $M$ (\hyperref[II.8.2.5-fr]{8.2.5}), cette image
s'identifie à $N$, en vertu de la définition de $\overline{N}$ donnée dans
(\hyperref[II.8.13.2-fr]{8.13.2}).

Pour prouver le second point, notons que l'injection
$i:C\cap\widehat{C}_f\to\widehat{C}_f$ correspond à l'injection canonique
$S_f^{\leq}\to S_f$ (\hyperref[II.8.3.2.6-fr]{8.3.2.6})~; d'autre part, on
a $\Gamma(\widehat{C}_f,\mathcal{M}^{\square})=M_f^{\leq}$,
$\Gamma(\widehat{C}_f,i_*(\widetilde{\mathcal{N}}))=N_f$ et (par
(\hyperref[II.8.12.2.1-fr]{8.12.2.1}))
$\Gamma(\widehat{C}_f,i_*(i^*(\mathcal{M}^{\square})))=M_f$. Compte tenu de
(\hyperref[I.9.4.2-fr]{I, 9.4.2}), on est donc ramené à montrer que
$\overline{N}_{(f)}\subset\widehat{M}_{(f)}=M_f^{\leq}$ s'identifie
canoniquement à l'image réciproque de $N_f$, par l'injection canonique
$M_f^{\leq}\to M_f$. En effet, soit $d=\deg(f)>0$, et supposons qu'un
élément $(\sum_{k\leqslant md}x_k)/f^m$ de $M_f$ (avec $x_k\in M_k$) soit de
la forme $y/f^m$ avec $y\in N$. En multipliant $y$ et les $x_k$ par un même
$f^h$ convenable, on peut déjà supposer que $\sum_{k\leqslant md}x_k=y$.
Mais dans l'identification (\hyperref[II.8.2.5.2-fr]{8.2.5.2}),
$(\sum_{k\leqslant md}x_k)/f^m$ correspond à
$\sum_{k\leqslant md}x_kz^{md-k}/f^m$, et cet élément appartient bien à
$\overline{N}_{(f)}$, puisque $\sum_{k\leqslant md}x_k\in N$~; la réciproque
est évidente.
\end{proof}

\begin{remark}[8.13.4]
\label{II.8.13.4-fr}
\begin{enumerate}
  \item[\rm{(i)}] Le cas d'application le plus important de
  (\hyperref[II.8.13.3-fr]{8.13.3}) est celui où
  $\mathcal{M}=\mathcal{S}$, $\widetilde{\mathcal{N}}$ étant donc un faisceau
  quasi-cohérent \emph{arbitraire} d'idéaux $\mathcal{J}$ de
  $\mathcal{O}_C$ (\hyperref[II.1.4.3-fr]{1.4.3}), correspondant
  biunivoquement à un \emph{sous-préschéma fermé} $Z$ de $C$. Alors le
  prolongement canonique $\overline{\mathcal{J}}$ de $\mathcal{J}$ est le
  faisceau quasi-cohérent d'idéaux de $\mathcal{O}_{\widehat{C}}$ qui définit
  l'\emph{adhérence} $\overline{Z}$ de $Z$ dans $\widehat{C}$
  (\hyperref[I.9.5.10-fr]{I, 9.5.10})~; la prop.
  (\hyperref[II.8.13.3-fr]{8.13.3}) donne un moyen canonique de définir
  $\overline{Z}$ à l'aide d'un Idéal gradué dans
  $\widehat{\mathcal{S}}=\mathcal{S}[z]$.
  \item[\rm{(ii)}] Supposons pour simplifier que $Y$ soit affine, et gardons
  les notations de la démonstration de
  (\hyperref[II.8.13.3-fr]{8.13.3}). Pour tout $x\in N$ non nul, soit $d(x)$
  le plus grand des degrés des composantes homogènes $x_i$ de $x$ dans $M$~;
  $\overline{N}$ est par définition le sous-module de $\widehat{M}$ formé de
  $0$ et des éléments de la forme
  $h(x,k)=z^k\sum_{i\leqslant d(x)}x_i z^{d(x)-i}$ ($k$ entier
  $\geqslant0$)~; il est donc engendré, en tant que module sur
  $\widehat{S}=S[z]$, par les éléments de la forme
  \[
    h(x,0)=\sum_{i\leqslant d(x)}x_i z^{d(x)-i}.
  \]

  